% Replacement simulation methods and results section for wcf_main.tex.
% The values below are read from the final cwdb_dr adapter (labelled WCF),
% original-code DRF/Causal-DRF records, and the phase-6.5 zero-inflation run.

\section{Simulation studies}\label{sec:simulations}

We evaluate WCF in finite-grid distributional treatment-effect problems with
observed covariates, binary adoption, and a quantile vector recorded for each
unit. The simulation studies include abstract distributional regimes, an
income-shaped family with endogenous adoption, and a zero-inflated family for
which a point mass is part of the outcome law. All comparisons use the same
training sample, test design, and random seed within a replication. The
results reported here use the final WCF specification throughout.

\subsection{Data-generating process and implementation}

For each unit, let $X\in[-1,1]^p$, $A\in\{0,1\}$, and let $q(Y^a)$ denote a
quantile vector on the grid
$u_k=(k-1/2)/K$, with $z_k=\Phi^{-1}(u_k)$ and equal weights
$w_k=1/K$. The income-shaped family is generated by
\[
q_k(Y^a)=m_a(X)+\xi+\exp\{s_a(X)+\eta\}\psi(z_k;\gamma_a(X)),
\qquad
\psi(z;\gamma)=z+\gamma\frac{z^2-1}{2}
 +\gamma\frac{z^3-3z}{6}.
\]
The location shock and log-scale shock are independent of $(X,A)$ within an
arm. The four regimes are a placebo, an earned-income-credit-like location
and shape intervention, a wage-floor-like location and scale intervention,
and the earned-income-credit-like intervention under reduced overlap. Their
assignment probabilities are clipped logistic functions of the covariates.
The abstract family uses the same finite-grid construction and includes
deterministic, null, nonlinear, multimodal, shape-transfer, and strong-
confounding regimes. The zero-inflated family replaces the law by
\[
 L_a(x)=\{1-\pi_a(x)\}\delta_{0_K}+\pi_a(x)L_a^+(x),
\]
where $0_K$ is the all-zero quantile vector and $L_a^+$ is a positive-part
law.

WCF represents a conditional law by $M$ particles in the monotone cone and
fits the particles by gradient boosting against a smoothed energy score. One
shared covariate partition is used for both arms, with cross-validated
arm-contrast shrinkage. Its calibrated functional estimates use the
cross-fitted augmented inverse-propensity score
\[
\begin{split}
\varphi_i^{(h)}={}&\widehat\mu_1^{(h)}(X_i)-\widehat\mu_0^{(h)}(X_i)
 +\frac{A_i}{\widehat e(X_i)}\{h(q_i)-\widehat\mu_1^{(h)}(X_i)\}\\
&-\frac{1-A_i}{1-\widehat e(X_i)}
 \{h(q_i)-\widehat\mu_0^{(h)}(X_i)\}.
\end{split}
\]
Here $\widehat e$ is estimated by cross-fitted logistic regression and clipped
to $[0.02,0.98]$ for the reported simulations. Logistic regression is an
implementation choice for this study; the score itself permits any
cross-fitted binary probability estimator. The calibrated layer changes the
functional and reference contrasts but leaves the predicted particle laws
unchanged. The optional zero-inflated extension estimates $\pi_a(x)$ with a
per-arm penalized logistic classifier and fits the particle booster to the
positive component.

The primary specification uses $K=25$, $M=10$, 100 boosting steps, learning
rate $0.12$, tree depth 4, minimum leaf size 10, minimum arm-leaf size 5,
collision parameter $10^{-3}$, and candidate shrinkage strengths
$\{0,50,500\}$. These numerical choices are examined in a separate
resolution study described in Section~\ref{sec:sensitivity-plan}; they are
not treated as theoretically privileged.

\subsection{Estimands and evaluation}

The mean-quantile contrast is
\[
\tau_q(x)=\mathbb{E}\{q(Y^1)-q(Y^0)\mid X=x\},
\]
where the reported `meanQ' metric is the RMSE of this full $K$-vector. For a
grid functional $h$, the marginal and moderator-conditional distributional
effects are
\[
\operatorname{TATE}_h=\mathbb{E}\{h(q(Y^1))-h(q(Y^0))\},\qquad
\operatorname{TCATE}_h(v)=\mathbb{E}\{h(q(Y^1))-h(q(Y^0))\mid V=v\}.
\]
We use the grid mean, grid standard deviation, grid skewness, and upper-tail
mean as $h$. The reference-distribution effects are a second pair of
estimands:
\[
r_a^K(x)=\mathbb{E}\{d_{W,K}(q(Y^a),q_\star)\mid X=x\},
\]
\[
\operatorname{Ref\!-!ATE}=\mathbb{E}\{r_1^K(X)-r_0^K(X)\},\qquad
\operatorname{Ref\!-!TCATE}(v)=\mathbb{E}\{r_1^K(X)-r_0^K(X)\mid V=v\}.
\]
Thus the reference metrics quantify whether treatment moves the outcome law
toward or away from the declared benchmark distribution. They are the
reference-distribution analogues of TATE and TCATE.

For each replication, law metrics average the arm-specific errors over the
test design. TATE and Ref-ATE are absolute errors of the marginal contrast;
TCATE and Ref-TCATE are RMSEs over the four moderator bins. The tables report
the mean of the replication-level errors over seeds. The law metric is the
squared Gaussian-kernel MMD to the oracle conditional law, and lower values
are better for every metric.

\subsection{Results for income-shaped outcomes}

Table~\ref{tab:revision-income} reports the results at $n=1000$, averaged
over ten seeds. WCF has the smallest error for every listed target in the
four income regimes. The placebo IC0 is especially informative: all causal
effects are zero, yet the two forest baselines produce substantially larger
false mean, functional, and reference effects. Under IC3, the reduced-overlap
regime, WCF's reference-TCATE error remains below the two forest baselines.
These comparisons support the performance claim for the stated income-shaped
designs and overlap range.

\begin{table}[t]
\centering\scriptsize
\caption{Income-shaped simulation results at $n=1000$. Entries are means over
ten seed-level errors. TATE and TCATE average the four registered grid
functionals within a seed. Ref-ATE and Ref-TCATE are the two
reference-distribution effects.}
\label{tab:revision-income}
\begin{tabular}{llrrrrrr}
\toprule
Regime & Method & meanQ & Law & TATE & TCATE & Ref-ATE & Ref-TCATE\\
\midrule
IC0 & WCF & .0356 & .0153 & .0100 & .0240 & .0142 & .0298\\
    & Causal-DRF & .1947 & .0467 & .1134 & .1138 & .1570 & .1572\\
    & DRF & .1963 & .0641 & .1104 & .1105 & .1543 & .1544\\
IC1 & WCF & .0904 & .0174 & .0128 & .0271 & .0143 & .0300\\
    & Causal-DRF & .1760 & .0455 & .0962 & .0966 & .1329 & .1333\\
    & DRF & .1876 & .0591 & .1012 & .1013 & .1394 & .1395\\
IC2 & WCF & .0700 & .0167 & .0099 & .0255 & .0143 & .0330\\
    & Causal-DRF & .1975 & .0488 & .1176 & .1179 & .1565 & .1567\\
    & DRF & .2298 & .0724 & .1357 & .1358 & .1795 & .1796\\
IC3 & WCF & .1103 & .0180 & .0188 & .0412 & .0161 & .0463\\
    & Causal-DRF & .3177 & .0605 & .1784 & .1786 & .2474 & .2476\\
    & DRF & .3481 & .0773 & .1972 & .1973 & .2701 & .2702\\
\bottomrule
\end{tabular}
\end{table}

The abstract suite gives a more qualified picture. At $n=1000$, WCF has the
smallest law error in D0, D2, D5, D7, and D8, but not in D6, whose outer law
is multimodal. WCF also has larger Ref-TCATE error than both incumbents on
D6. The final paper should therefore state that WCF is effective on the
income-shaped family and on several abstract regimes, while its performance
depends on the geometry of the conditional law.

\begin{table}[t]
\centering\scriptsize
\caption{Abstract-suite results at $n=1000$. Entries are means over seeds
0--9. Incumbent rows use the original-code merged records with the same seed
restriction.}
\label{tab:revision-abstract}
\begin{tabular}{lrrrrrrrrr}
\toprule
 & \multicolumn{3}{c}{meanQ} & \multicolumn{3}{c}{Law} &
 \multicolumn{3}{c}{Ref-TCATE}\\
DGP & WCF & Causal-DRF & DRF & WCF & Causal-DRF & DRF &
 WCF & Causal-DRF & DRF\\
\midrule
D0 & .0868 & .1391 & .1672 & .0234 & .1033 & .1191 & .0097 & .0886 & .0799\\
D2 & .0410 & .0818 & .0627 & .0098 & .0154 & .0138 & .0399 & .0227 & .0224\\
D5 & .0946 & .0765 & .0917 & .0220 & .0374 & .0431 & .0418 & .0239 & .0235\\
D6 & .2189 & .2310 & .1986 & .0178 & .0064 & .0059 & .0862 & .0585 & .0501\\
D7 & .0553 & .0721 & .0662 & .0094 & .0129 & .0142 & .0386 & .0235 & .0238\\
D8 & .0635 & .2100 & .1623 & .0127 & .0417 & .0416 & .0793 & .0584 & .0724\\
\bottomrule
\end{tabular}
\end{table}

\subsection{Zero-inflated outcomes}

The ordinary particle representation cannot place positive probability on the
all-zero quantile vector. Table~\ref{tab:revision-zi} shows why the optional
two-part assembly is needed. At $n=1000$, WCF-ZIPT reduces zero-mass and law
error in every zero-inflated regime. WCF-ZIPT is reported as an uncalibrated
two-part extension: its classifier and positive-part booster do not include
the AIPW functional calibration used by WCF.

\begin{table}[t]
\centering\scriptsize
\caption{Zero-inflated simulation results at $n=1000$. Zero-mass is the RMSE
of implied mass at the degenerate component, mass contrast is the moderator-
binned error in the treatment contrast of that mass, and Law is squared MMD.}
\label{tab:revision-zi}
\begin{tabular}{llrrr}
\toprule
Regime & Method & Zero-mass & Mass contrast & Law\\
\midrule
ZI0 & WCF & .4947 & .0094 & .1113\\
    & WCF-ZIPT & .0537 & .0427 & .0133\\
    & Causal-DRF & .0945 & .0595 & .0247\\
    & DRF & .1058 & .0561 & .0296\\
ZI1 & WCF & .4173 & .1532 & .1264\\
    & WCF-ZIPT & .0526 & .0365 & .0154\\
    & Causal-DRF & .0901 & .0426 & .0249\\
    & DRF & .1022 & .0479 & .0291\\
ZI2 & WCF & .4954 & .0080 & .1112\\
    & WCF-ZIPT & .0537 & .0427 & .0136\\
    & Causal-DRF & .0946 & .0597 & .0258\\
    & DRF & .1060 & .0564 & .0304\\
ZI3 & WCF & .4951 & .0746 & .1741\\
    & WCF-ZIPT & .0598 & .0445 & .0176\\
    & Causal-DRF & .1731 & .2051 & .0768\\
    & DRF & .3475 & .3823 & .2472\\
\bottomrule
\end{tabular}
\end{table}

The (K)- and (M)-sensitivity study and the symmetric nonlinear-assignment
study are specified in Section~\ref{sec:sensitivity-plan}. Until those cells
are run, the paper should present $K=25$ and $M=10$ as the primary
implementation choices for these simulations rather than as resolution-free
defaults.
